% Einheit 8 von 8 – Rückwärts (bank/zufallsexperimente-und-pfadregeln/e8.jsonl, zusammenbau v0.1)
%% TODO Zeitmarke Einheit 8: Marken-Zeile nicht lesbar
%% TODO Zweigzeile Teil 1: was hier gelernt wird (Satz in Schülersprache aus dem Einheitstitel) – Einheit 8
\einheitenkopf[e8]{Einheit 8 von 8 · Rückwärts}
\zweigzeile{Abitur GK}
\setcounter{aufgabe}{28}

%% TODO Titel: Ich-kann-Satz fehlt in der Bank (Platzhalter: Kettenname) – Nr. 29
%% TODO Anweisung: ein Satz über der ersten Teilaufgabe fehlt in der Bank – Nr. 29
\begin{aufgabe}{Rückwärts}
%% TODO \rechenplatz{n}: Schrittzahl des Grundfalls fehlt in der Bank (nur bei fester Schreibform, 2.3 b)
\begin{teile}
\teil „Eine Urne enthält 3 rote und 5 blaue Kugeln. Wie wahrscheinlich sind zwei rote beim Ziehen ohne Zurücklegen?“\\ \kreuz{vorwärts: eine Wahrscheinlichkeit ist gesucht}\\ \kreuz{rückwärts: eine Größe im Baum ist gesucht}
\end{teile}
\begin{gleichungsraster}[2]
\gl{\text{60\,\% der Kunden sind Frauen, unter ihnen sind 50\,\% Stammkunden. Insgesamt sind 44\,\% Stammkunden. Anteil a der Stammkunden unter den Männern?}} &
\gl{\text{Befragung mit Zufallsfrage: Mit $0{,}7$ kommt F1 „Hast du schon einmal …?“, mit $0{,}3$ F2 „Hast du noch nie …?“. 50\,\% antworten mit Ja. Anteil p derer, die es schon getan haben?}} \\
\gl{\text{In einem Betrieb arbeiten 5\,\% der Frauen und 15\,\% der Männer im Schichtdienst. Der Anteil aller Beschäftigten, die Männer im Schichtdienst sind, ist doppelt so groß wie der Anteil der Frauen im Schichtdienst. Frauenanteil w?}} &
\gl{\text{Ein Glücksrad hat einen roten und einen weißen Sektor; es wird zweimal gedreht. P(zweimal dieselbe Farbe) $= 0{,}58$. Mittelpunktswinkel des kleineren Sektors?}} \\
\gl{\text{Eine Urne enthält n Kugeln, 3 davon rot. Zwei werden ohne Zurücklegen gezogen; P(beide rot) $= \frac{1}{15}$. n?}} &
\gl{\text{Wie oft muss man mindestens würfeln, damit P(mindestens eine Sechs) wenigstens $0{,}9$ ist?}} \\
\gl{\text{Ein Glücksrad hat einen grünen Sektor mit P(grün) $= 0{,}4$, dazu einen roten und einen gelben. Es wird zweimal gedreht. A: einmal Rot und einmal Gelb. Der rote Sektor ist so gewählt, dass P(A) maximal ist. Mittelpunktswinkel des roten Sektors? (Abitur 2022 GK)}} & \\
\end{gleichungsraster}
\end{aufgabe}

%% TODO Titel: Ich-kann-Satz fehlt in der Bank (Platzhalter: Kettenname) – Nr. 30
%% TODO Anweisung: ein Satz über der ersten Teilaufgabe fehlt in der Bank – Nr. 30
\begin{aufgabe}{Rückwärts – Fehler finden, Begründen, Anwendung, Darstellungswechsel}
\begin{teile}
\teil Zufallsfrage: F1 mit $0{,}7$, F2 („noch nie …“) mit $0{,}3$; 42\,\% sagen Ja. Leo: „Also ist $p = 0{,}42$.“ Finde den Fehler.
\teil Begründe: Eine gegebene Randwahrscheinlichkeit liefert eine Gleichung für einen unbekannten Ast.
\end{teile}
\begin{gleichungsraster}[2]
\gl{\text{In einer Fahrschule bestehen 75\,\% die Prüfung beim ersten Mal; von den übrigen besteht der Anteil a beim zweiten Mal. Insgesamt bestehen 93\,\% spätestens beim zweiten Mal. a?}} &
\gl{\text{30\,\% gehören zu Gruppe A, von ihnen haben 80\,\% ein Merkmal; in Gruppe B ist der Anteil x unbekannt. Insgesamt haben 62\,\% das Merkmal. Stelle die Gleichung auf und löse sie.}} \\
\end{gleichungsraster}
\end{aufgabe}
